\documentclass[12pt,twosides,onecolumn,openany]{article}
\usepackage{graphicx} 

\usepackage{cite}
\usepackage[catalan]{babel}
\usepackage{emptypage}
\usepackage{hyperref}
\usepackage{mathtools}
\usepackage{blindtext}
\usepackage{subfigure}
\usepackage[utf8]{inputenc}
\usepackage{caption}
\usepackage{subcaption}
\usepackage{wrapfig}
\usepackage[a4paper]{geometry}
\geometry{top=2.5cm, bottom=2.5cm, left=2.5cm, right=2.5cm}
\usepackage{fancyhdr}
\pagestyle{fancy}
\usepackage{amsmath}
\usepackage{amssymb}
\usepackage{amsfonts}
\providecommand{\norm}[1]{\lVert#1\rVert}
\hypersetup{colorlinks=true,urlcolor=blue,linkcolor=blue}
\usepackage{multirow}
\usepackage{multicol}
\usepackage{rotating}

\usepackage{titlesec}

\newenvironment{Figura}
  {\par\medskip\noindent\minipage{\linewidth}}
  {\endminipage\par\medskip}

\titleformat{\section}  % comando de sección a formatear
  {\fontsize{14}{16}\bfseries} % formato para toda la línea
  {\thesection} % cómo mostrar el número
  {0.4em} % espacio entre el número y el texto
  {} % formato solo para el texto
  [] % formato para después del texto


\fancyhf{}
\fancyhead[LO,RE]{Grup D3}
\fancyhead[RO,LE]{Pràctica Ia}
\fancyfoot[LO,RE]{\thepage}
\fancyfoot[RO,LE]{Laboratori de Termodinàmica - UAB}
\graphicspath{ {images/} }

\begin{document}

\begin{center}
    {\Large \textsc{Transport de la calor: Estudi del transport de la calor en una barra metàl·lica}}\\
    \vspace{0.2cm}
    \textsc{30 d'Octubre de 2024}\\
    \vspace{0.2cm}
    $\begin{matrix} 
    \text{Miguel A.} \hspace{1.5cm} & \text{Daniel B.} & \hspace{1.5cm} \text{Sergi R.}\\1637738 \hspace{1.5cm} & 1603508 & \hspace{1.5cm} 1607805
    \end{matrix}$
\end{center}
\begin{center}
    \textsc{\textit{RESUM}}
\end{center}
En aquesta pràctica s'estudien els fenòmens de transport de la calor en metalls. Es construeix un model teòric del nostre sistema a partir de les equacions de Fourier i Newton i de la llei d'Stefan-Boltmann. A partir de dos sistemes formats principalment per barres metàl·liques i una manta calefactora comprovem experimentalment la coherència del model teòric presentat i discutim les seves limitacions i assumpcions. 
\vspace{0.5cm}
\begin{multicols}{2}
\section{Introducció i objectius}
El transport de calor es pot realitzar de tres maneres diferents: conducció, convecció i radiació. Totes tres formes participen alhora en qualsevol sistema termodinàmic clàssic.\\\\
La conducció d'una placa d'espesor infinitesimal ve descrita per la llei de Fourier:
\begin{equation}\label{Fourier}
  \frac{\text{d}Q}{\text{d}t} = -KA \frac{\text{d}T}{\text{d}x}
\end{equation}
on $K$ és la conductivitat tèrmica i $A$ l'àrea la placa. La conducció considera el transport de calor essencialment de la matèria condensada, on dos sistemes en contacte intercanvien energia sense haver transport de matèria.\\\\
La convecció és un fenòmen que té en compte l'intercanvi d'energia per mitjà de l'intercanvi de matèria. Tot i ser complicat de desenvolupar-la teòricament, com a aproximació d'un sistema on intervé un sòlid i un fluid, podem descriure la convecció amb la llei empírica de refredament de Newton:
\begin{equation}\label{Newton}
  \frac{\text{d}Q}{\text{d}t} = -\lambda A(T-T_a)
\end{equation}
on $\lambda$ és el coeficient de convecció, $A$ l'àrea de contacte entre el sòlid i el fluid i $T_a$ la temperatura del fluid.\\\\
La radiació és l'intercanvi de calor sense necessitat de cap material, només a través d'ones electromagnètiques. Aquest mecanisme ve descrit per la llei d'Stefan-Boltzmann:
\begin{equation}\label{Stefan-Boltzmann}
  \frac{\text{d}Q}{\text{d}t} = e\sigma_{\text{SB}}AT^{4}
\end{equation}
on $e$ és l'emisivitat característica del cos que conforma el sistema, \(\sigma_{\text{SB}} = (5.6697\pm0,0029) \cdot10^{-8}\,\text{W}/\text{m}^{2}\text{K}^{4}\) és la constant d'Stefan-Boltzmann i $A$ és la superfície del cos.\\\\
Durant aquesta pràctica comprovarem com i quant aporten aquestes contribucions de calor per a dos sistemes diferents, formats per barres de metalls diferents. Considerant una barra sòlida amb una forma cilíndrica ideal podem construïr un model teòric per al nostre sistema. Menyspreant les contribucions per rafiació (més endavant comprovarem que això no causa problemes) arribem a l'equació diferencial
\begin{equation}\label{EDP}
  \frac{\partial \theta_x}{\partial x^2} = \frac{2\lambda}{Kr}\theta_x + \frac{c\rho}{K}\frac{\partial \theta_x}{\partial t}
\end{equation}
on $\theta_x$ és la diferència de temperatura respecta a l'ambient a la posició $x$ i $c$ i $\rho$ la capacitat calorífica i la densitat del material, respectivament.\\\\ Considerant un sistema estacionari per a resoldre l'Equació \eqref{EDP} obtenim la solució:
\begin{equation}\label{sol_estacionaria}
  \theta(x) = \theta_0e^{-px}
\end{equation}
on $p = \sqrt{2\lambda/Kr}$. Considerant la temperatura a l'origen una funció periòdica obtenim la solució:
\begin{multline}\label{sol_permanent}
  \theta(x,t) = \theta_0e^{-px} +\\
   + A_0e^{-mx}\cos{\left( \frac{2\pi}{\tau} -hx \right)} 
\end{multline}
on $\tau$ és el periode d'oscil·lació i es compleix $p^2 = m^2-h^2$ i $mh = \rho \pi c / K \tau$. Si fem una mitjana temporal d'aquesta solució recuperem la solució estacionaria:
\begin{equation}\label{prom_temp}
  \langle \theta(x,t) \rangle = \theta_0e^{-px}  
\end{equation}
Si ens fixem novament en la solució permanent, podem observar que eisteix un desfasament de les ones al llarg de l'espai
\begin{equation}\label{desfasament}
  \varphi_i = hx_i
\end{equation}
aquest fet serà útil per a calcular la variable $h$ ja que també es complirà que
\begin{equation}\label{increment_desfasament}
  \Delta \varphi_i = h\Delta x
\end{equation}
Notem que les oscil·lacions tenen una amplitud que decau amb la distància, sabent això podem calcular la variable $m$ fent
\begin{equation}\label{trobar_m}
  \frac{a_i}{a_j} = e^{-m(x_i-x_j)}
\end{equation}
Gracies a aquestes definicions, podem trobar també un valor per al coeficient de convecció de forma que
\begin{equation}\label{valor_lambda}
  \lambda = \frac{1}{2}[Kr(m^2 - h^2)]
\end{equation}
El nostre sistema estacionari està conformat per tres barres metàl·liques \cite{guia_lab}. Una és d'alumini (Al), una altra de llautó (Cu-Zn) i una altra de ferro (Fe). Totes tres de la mateixa longitud i de radis similars. A un extrem de cada barra situem una manta calefactora que els mantindrà a una temperatura elevada. Per a evitar que l'escalfament de l'aire per conducció afecti a les aportacions per convecció del nostre sistema fem servir un ventilador que mou l'aire al voltant de l'extrem i del mig de les barres. Cada barra té un conjunt de forats situats a la posició $x_i$ respecte al forat més a prop de la manta calefactora, separats entre ells per una distància \(d = (10,\pm0.5)\, \text{cm}\). Les mesures de temperatura les realitzem de des de l'extrem fred fins a l'extrem calent, passant pel forat $x_i$ de les tres barres consecutivament. L'objectiu d'això és aminorar el temps de la presa de mesures, ja que en cas contrari hauriem d'esperar entre mesura i mesura a que el termòmetre torni a rondar la temperatura ambient.\\\\
El sistema permanent consisteix de dues barres d'alumini, de gruix diferent cadascuna, connectades de forma similar al sistema estacionari. La manta calefactora està connectada a un sistema automatitzat que escalfa les barres de forma periòdica. També en recull les dades de temperatura provinents de tres termoparells situats a les posicions (respecte al primer termoparell) \(0\,\text{cm}\), \(10\,\text{cm}\) i \(15\,\text{cm}\) per a la barra gran i \(0\,\text{cm}\), \(10\,\text{cm}\) i \(20\,\text{cm}\) per a la barra petita. 
\section{Resultats i discussió}
\subsection{Estat estacionari}
Segons el desenvolupament teòric que seguim no podem negligir les diferències de radi entre les tres barres metàl·liques. Mesurem a l'inici, al mig i al final el diàmetre de cada barra i obtenim els valors promig de la Taula \ref{Tau:diam_radis}.
\begin{Figura}
  \centering
  \captionof{table}{\footnotesize{Valors promig del radi de cada barra}}
  \begin{tabular}{c|c}
    Material & $r$ ($\pm$0,001) [cm] \\
    \hline\hline
    Al & 1,524\\
    Cu-Zn & 1,518\\
    Fe & 1,573 
  \end{tabular}
  \label{Tau:diam_radis}
\end{Figura}
A més, al llarg de l'experiment el labortaori ha anat canviant lleugerament de temperatura ambient. Per aquest motiu a l'hora de calcular els increments de temperatura $\theta_x$ farem servir com a temperatura ambient de referència la mitjana de tres mesures (veure Taula \ref{T_ambient} a l'Annex). Per tant, el valor de la temperatura ambient al laboratori és: \(T_a = (24,73 \pm 0,19)\,^\circ\text{C}\).\\\\

A partir de les dades de temperatura obtingudes per a cada barra obtenim la gràfica de la Figura \ref{DeltaT_vs_d}. Es pot veure com l'increment de temperatura decau de forma exponencial amb la distància al punt calent, tal com es podia preveure del desenvolupament teòric. Fent servir codi en Python obtenim un ajust exponencial de les dades mesurades. Aquest ajust es pot observar també a la Figura \ref{DeltaT_vs_d}, i les dades dels coeficients ajustats a la Taula \ref{Tau:coef_exponencial}.
\begin{Figura}
  \centering
  \captionof{table}{\footnotesize{Coeficients $a$ i $b$ de l'ajust exponencial, amb model $y = ae^{bx}$, per als punts experimentals mesurats per l'alumini, el llautó i el ferro.}}
  \begin{tabular}{c|c|c}
    Material & $a$ [$^\circ$C] & $b$ ($\cdot 10^{-2}$) [cm$^{-1}$]\\
    \hline\hline
    Al & $95,11\pm0,97$ & $-3,025\pm0,051$\\
    Cu-Zn & $81,5\pm2,2$ & $-4,53\pm0,21$ \\
    Fe & $95,1\pm1,3$ & $-8,03\pm0,22$
  \end{tabular}
  \label{Tau:coef_exponencial}
\end{Figura}

\begin{Figura}
  \centering
  \includegraphics[width = 1\linewidth]{../../graphs/practica_Ia/plots/theta_vs_d_estacionaria.png}\label{DeltaT_vs_d}
  \captionof{figure}{\footnotesize{Increment de temperatura $\theta$ respecte a la disància, punts experimentals i regressió exponencial. Coeficients de regressió: $r^2_{\text{Al}} = 0.999$, $r^2_{\text{Cu-Zn}} = 0.991$ i $r^2_{\text{Fe}} = 0.998$}}
\end{Figura}

Per a trobar el valor experimental de $p$ de l'Equació \eqref{prom_temp} apliquem logaritme als valors experimentals per a poder realitzar una regressió lineal. Observem aquestes dades experimentals i les rectes de regressió a la Figura \ref{lnDeltaT_vs_d}. A la Taula \ref{Tau:dades_regressio} s'observen els valors ajustats de les regressions.
\begin{Figura}
  \centering
  \captionof{table}{\footnotesize{Coeficients $a$ i $b$ de l'ajust lineal, amb model $y = ax + b$, per al logaritme de l'increment de temperatura $\theta$ per a l'alumini, el llautó i el ferro.}}
  \begin{tabular}{c|c|c}
    Material & $a$ ($\cdot 10^{-2}$) [cm$^{-1}$] & $b$\\
    \hline\hline
    Al & $-2,97\pm0,12$ & $4,517\pm0,080$\\
    Cu-Zn & $-3,40\pm0,11$ & $4,164\pm0,070$ \\
    Fe & $-5,20\pm0,26$ & $4,07\pm0,17$
  \end{tabular}
  \label{Tau:dades_regressio}
\end{Figura}

\begin{Figura}
  \centering
  \includegraphics[width = 1\linewidth]{../../graphs/practica_Ia/plots/reg_estacionaria.png}\label{lnDeltaT_vs_d}
  \captionof{figure}{\footnotesize{Regressió lineal del logaritme de les temperatures per a cada material. Coeficients de regressió obtinguts: $r^2_{\text{Al}}=0.983$, $r^2_{\text{Cu-Zn}}=0.990$, $r^2_{\text{Fe}}=0.976$.}}
\end{Figura}           
Tenint en compte de l'equació \eqref{sol_estacionaria}, què val $p$ (el valor absolut de les $a$ a la Taula \ref{Tau:dades_regressio}) i suposant el mateix valor constant del coeficient $\lambda$ per als tres metalls observem la temperatura d'un metall decau com $(Kr)^{1/2}$. Deduim doncs que la conductivitat tèrmica dels tres metalls és diferent i que es compleix
\begin{equation*}
  K_{\text{Fe}} < K_{\text{Cu-Zn}} < K_{\text{Al}} \hspace{3mm}
\end{equation*}
i que existeixen les relacions
\begin{equation*}
  \frac{K_ir_i}{K_jr_j} = \frac{p^2_j}{p^2_i}
\end{equation*}
on $i$ i $j$ prenen valors dels diferents metalls que estem estudiant (Al, Cu-Zn i Fe). Fent recerca a la literatura trobem el valor de la conductivitat tèrmica del ferro\footnote{Considerarem aquestes dades com a patró o referència, per tant no cal que els hi apliquem una incertesa.}\cite{CRC}:
\begin{equation*}
  K_{\text{Fe}} =  0,802\, \text{W}/\text{cmK}
\end{equation*}
Fent servir les relacions anteriors i el valor trobat a la literatura podem calcular, aproximadament, els valors de la conductivitat tèrmica dels altres dos metalls. Podem comparar aquests valors obtinguts amb els que podem trobar a la literatura. S'observen totes aquestes dades a la Taula \ref{Tau:conductivitats_termiques_Fe}.
\begin{Figura}
  \centering
  \captionof{table}{\footnotesize{Comparació de la conductivitat tèrmica experimental i teòrica segons el patró ferro.}\cite{CRC}\cite{Young}}
  \begin{tabular}{c|c|c}
    Material & $K_{\text{exp}}$ [W/cmK] & $K_{\text{teo}}$ [W/cmK]\\
    \hline\hline
    Al & $2,54\pm0,28$ & 2,37\\
    Cu-Zn & $1,94\pm0,21$ & 1,09
  \end{tabular}
  \label{Tau:conductivitats_termiques_Fe}
\end{Figura}
Com a aproximació alternativa podem fer servir, en comptes del ferro, l'alumini com a patró. Aplicant les mateixes relacions que amb el ferro obtenim les dades de la Taula \ref{Tau:conductivitats_termiques_Al}.
\begin{Figura}
  \centering
  \captionof{table}{\footnotesize{Comparació de la conductivitat tèrmica experimental i teòrica segons el patró alumini.}\cite{CRC}\cite{Young}}
  \begin{tabular}{c|c|c}
    Material & $K_{\text{exp}}$ [W/cmK] & $K_{\text{teo}}$ [W/cmK]\\
    \hline\hline
    Cu-Zn & $1,81\pm0,17$ & 1,09\\
    Fe & $0,748\pm0,078$ & 0,802
  \end{tabular}
  \label{Tau:conductivitats_termiques_Al}
\end{Figura}
Observant aquestes dues darreres taules notem com el càlcul experimental de $K_{\text{Al}}$ segons el patró ferro i de $K_{\text{Fe}}$ segons el patró alumini són coherents amb les dades obtingudes de la literatura d'aquests mateixos coeficients. En canvi, el valor de $K_{\text{Cu-Zn}}$ tant per al patró ferro com per al patró alumini difereix, el valor que trobem a la literatura no entra dins de les incerteses. Això ens porta a deduïr que possiblement la hipòtesi inicial de considerar les $\lambda$ iguals per a tots tres materials no era correcta. També cal considerar que la constant de conductivitat tèrmica pot dependre de la temperatura i d'altres propietats micorscòpiques de cada material, i considerar-la constant és una simplificació.\cite{Ausloos}\\\\
D'igual manera, les pèrdues per convecció sofreixen una altra simplificació a l'hora de calcular l'àrea de contacte amb l'aire, el nostre model considera una barra sense forats. Tot i això aquesta simplificació no ens porta a obtenir resultats incoherents amb la teoria. Si les barres fossin totalment buides aquest canvi de l'àrea seria rellevant i canviarien els nostres resultats experimentals, hauriem de revisar el desenvolupament teòric.\\\\
Com a darrera simplificació que pateixen els nostres resultats podem fixar-nos en que no considerem els efectectes de la radiació. Comparem les aportacions de conductivitat, convecció i radiació. Observant les expressions de les Equacions \eqref{Fourier}, \eqref{Newton} i \eqref{Stefan-Boltzmann} podem observar que per al nostre sistema es compleixen els següents llindars:
\begin{equation*}
  \left|{P_{\text{F}}}\right| \leq \left| P_{\text{F}}(x=0) \right| 
\end{equation*}
per a l'aportació de l'equació de Fourier,
\begin{equation*}
  \left|P_{\text{N}} \right| \leq \left| P_{\text{N}}(T=T_{\text{max}}) \right|
\end{equation*}
per a l'aportació de l'equació de Newton i
\begin{equation*}
  \left|P_{\text{SB}} \right| \leq \left| P_{\text{SB}}(T=T_{\text{max}}) \right|
\end{equation*}
per a l'aportació de la llei d'Stefan-Boltzmann. Calculant aquests llindars tenint en compte que $T = \theta(x,t) + T_a$ obtenim els resultats de la Taula \ref{tau:comp_potencies}.
\begin{Figura}
  \centering
  \captionof{table}{\footnotesize{Taula de potència llindar del sistema per a cada material i cada aportació de calor. S'ometen les incerteses per a no distreure la comparació.}}
  \begin{tabular}{c|c|c|c}
    Material & $P_{\text{F}}$ [W] & $P_{\text{N}}$ [W] & $P_{\text{SB}}$ [W]\\ \hline\hline
    Al & 43,73 & 4810,66 & 0,004\\
    Cu-Zn & 17,74 & 1951,38 & 0,003\\
    Fe & 1,70 & 187,08 & 0,002
  \end{tabular}
  \label{tau:comp_potencies}
\end{Figura}
Cal aclarir novament que aquestes dades es tracten de llindars, la potència absorbida o emesa pel sistema és diferent a cada punt i depèn d'altres factors, per exemple del ventilador que manté un flux d'aire al voltant de les barres. Tot i això, les dades ens són suficients per a notar que l'aportació de pèrdues per radiació, seguint la llei d'Stefan-Boltzmann, és molt petita en comparació amb les altres dos. Sabent això concluim que podem menysprear del nostre model de l'aportació de la radiació.
\subsection{Estat permanent}
A partir de les dades obtingudes experimentalment representem l'evolució temporal de la temperatura segons cada posició de la barra gran, Figura \ref{fig:T_vs_t_gran}, i de la barra petita, \ref{fig:T_vs_t_petita}.
\begin{Figura}
  \centering
  \includegraphics[width= 1\linewidth]{../../graphs/practica_Ia/plots/gran.png}
  \captionof{figure}{\footnotesize{Gràfica de $T$ en funció $t$ de la barra gran. Temperatures mesurades a una distància de \(0\,\text{cm}\) (vermell), \(10\,\text{cm}\) (verd) i \(15\,\text{cm}\) (blau) del punt de referència.}}
  \label{fig:T_vs_t_gran}
\end{Figura}
\begin{Figura}
  \centering
  \includegraphics[width= 1\linewidth]{../../graphs/practica_Ia/plots/petita.png}
  \captionof{figure}{\footnotesize{Gràfica de $T$ en funció $t$ de la barra petita. Temperatures mesurades a una distància de \(0\,\text{cm}\) (vermell), \(10\,\text{cm}\) (verd) i \(20\,\text{cm}\) (blau) del punt de referència.}}
  \label{fig:T_vs_t_petita}
\end{Figura}
Notem com l'evolució temporal mesurada de la temperatura és coherent amb la solució de l'Equació \eqref{sol_permanent}. El segón terme de la solució correspón a una funció periòdica en el temps que decau amb la distància i presenta un desfasament, també depenent de la distància. Tant a la Figura \ref{fig:T_vs_t_gran} com a la Figura \ref{fig:T_vs_t_petita} es veu que la temperatura és funció periòdica en el temps (és una ona tèrmica) i que, quant més lluny del punt de referència ens trobem, més petita és l'amplitud de d'aquesta ona. Tot i ser complicat a simple vista veure-ho, existeix un desfasament entre les oscil·lacions de la mateixa barra a diferents distàncies. Per a observar millor aquest desfasament podem preparar les dades d'una forma més convenient.\\\\
Com a pas previ a observar el desfasament, calculem la temperatura mitjana a cada punt de cada barra. Com a simplificació, calculem aquesta mitjana temporal fent una mitjana aritmètica de les dades obtingudes al laboratori. Els valors d'aquestes mitjanes venen recollits a la Taula \ref{Tau:mitjanes}.
\begin{Figura}
  \centering
  \captionof{table}{\footnotesize{Temperatura mitjana ($\langle \theta \rangle$) a cada posició $x_i$ per a la barra gran i la barra petita.}}
  \begin{tabular}{c|c}
    $x_i$ ($\pm$0,1) [cm] & $\langle \theta \rangle$ ($\pm$ 0,001) [$^\circ$C]\\
    \hline 
    \multicolumn{2}{c}{Barra gran} \\ \hline
    0 & 81,849\\
    10 & 72,012\\
    15 & 67,928 \\ \hline
    \multicolumn{2}{c}{Barra petita} \\ \hline
    0 & 98,717\\
    10 & 80,420\\
    20 & 66,486 \\
  \end{tabular}
  \label{Tau:mitjanes}
\end{Figura}
Observem com la temperatura mitjana decau amb la distància respecte a l'origen de referència. Clarament són pocs punts experimentals com per a traçar una corba de tendència. Tot i això, qualitativament els valors obtinguts s'apropen al que prediu la teoria. En cas de repetir l'experiment amb més termoparells repartits al llarg de cadascuna de les barres observariem un comportament de la mitjana temporal de la temperatura com el de l'Equació \eqref{prom_temp}. Si apliquem el mateix mètode que a l'estat estacionari per a obtenir el valor de $p$ obtenim els valors que s'observen a la Taula \ref{tau:pendent_mitjana}. A partir d'aquestes dades realitzem una regressió i observem, Figura \ref{fig:lin_reg} com s'apropen molt a un comportament lineal.
\begin{Figura}
  \centering
  \captionof{table}{\footnotesize{Valors ajustats per a una regressió lineal, amb model $y = ax + b$, on $y$ és el logaritme de les temperatures mitjanes i $x$ les posicions $x_i$.}}
  \begin{tabular}{c|c|c}
    Barra & $a$ ($\cdot 10^{-3}$) [cm$^{-1}$] & $b$ ($\cdot 10^{-2}$) \\ \hline\hline
    gran & $-12,43\pm0,29$ & $44,04\pm0,30$\\
    petita & $-19,75\pm0,42$ & $45,90\pm0,55$
  \end{tabular}
  \label{Tau:pendent_mitjana}
\end{Figura}
\begin{Figura}
  \centering
  \includegraphics[width=1\linewidth]{../../graphs/practica_Ia/plots/linear_reg.png}
  \captionof{figure}{\footnotesize{Regressió lineal del logaritme de temperatures respecte a la posició. Coeficients de regressió obtinguts: $r^2$ = 0.999 tant per a la barra gran com per a la barra petita.}}
  \label{fig:lin_reg}
\end{Figura}
Un cop calculades les mitjanes per a cada punt, podem normalitzar les dades experimentals dividint per la seva temperatura mitjana corresponent. Aplicant aquest mètode i, per a obtenir més detall, escollint com a molt els dos primers periodes d'oscil·lació obtenim la Figura \ref{fig:T_vs_t_gran_norm} per a la barra gran i la Figura \ref{fig:T_vs_t_petita_norm} per a la barra petita.
\begin{Figura}
  \centering
  \includegraphics[width=1\linewidth]{../../graphs/practica_Ia/plots/gran_norm.png}
  \captionof{figure}{\footnotesize{Gràfica de $T$ en funció de $t$ normalitzada de la barra gran. Temperatures mesurades a una distància de \(0\,\text{cm}\) (vermell), \(10\,\text{cm}\) (verd) i \(15\,\text{cm}\) (blau) del punt de referència.}}
  \label{fig:T_vs_t_gran_norm}
\end{Figura}
\begin{Figura}
  \centering
  \includegraphics[width=1\linewidth]{../../graphs/practica_Ia/plots/petit_norm.png}
  \captionof{figure}{\footnotesize{Gràfica de $T$ en funció de $t$ normalitzada de la barra petita. Temperatures mesurades a una distància de \(0\,\text{cm}\) (vermell), \(10\,\text{cm}\) (verd) i \(20\,\text{cm}\) (blau) del punt de referència.}}
  \label{fig:T_vs_t_petita_norm}
\end{Figura}
Es pot obeservar com existeix desfasament entre dos punts diferents a cada barra, la ona tèrmica avança a mesura que augmenta la distància. A més, també es pot veure un comportament similar al que prediu l'Equació \eqref{desfasament}.\\\\
Per a determinar l'amplitud de la ona a cada punt per a cada barra escollim un punt màxim i un mínim consecutius i calculem la meitat de la diferència. Amb aquestes dades realitzem una mitjana i obtenim les amplituds que s'observen a la Taula \ref{tau:amplituds_mitjanes_gran} i a la Taula \ref{tau:amplituds_mitjanes_petita}.
\begin{Figura}
  \centering
  \captionof{table}{\footnotesize{Amplitud promig de la ona tèrmica sesgons la distància al punt de referència per a la barra gran.}}
  \begin{tabular}{c|c}
    $x_i$ ($\pm$0.1) [cm] & $\bar{a}_i$ ($\pm$0.001) [$^\circ$C]  \\ \hline
    0 & 2,963\\
    10 & 1,218\\
    15 & 0,802 
  \end{tabular}
  \label{tau:amplituds_mitjanes_gran}
\end{Figura} 
\begin{Figura}
  \centering
  \captionof{table}{\footnotesize{Amplitud promig de la ona tèrmica sesgons la distància al punt de referència per a la barra petita.}}
  \begin{tabular}{c|c}
    $x_i$ ($\pm$0.1) [cm] & $\bar{a}_i$ ($\pm$0.001) [$^\circ$C]  \\ \hline
    0 & 5,404\\
    10 & 2,273\\
    20 & 0,961  
  \end{tabular}
  \label{tau:amplituds_mitjanes_gran}
\end{Figura} 
Es pot observar com segueixen un comportament exponencial, similar al que hem vist a la Taula \ref{Tau:mitjanes}. Apliquem novament el mateix mètode amb les amplituds que amb les temperatures mitjanes i obtenim les regressions de la Figura \ref{fig:reg_lin_amplituds} i els valors dels pendents de la Taula \ref{tau:pendent_amplituds}.
\begin{Figura}
  \centering
  \captionof{table}{\footnotesize{Valors ajustats per a una regressió lineal, amb model $y = ax + b$, on $y$ és el logaritme de les amplituds i $x$ les posicions $x_i$.}}
  \begin{tabular}{c|c|c}
    Barra & $a$ ($\cdot 10^{-2}$) [cm$^{-1}$] & $b$ ($\cdot10^{-1}$) \\ \hline\hline
    gran & $-8,74\pm0,16$ & $10,82\pm0,14$\\
    petita & $-8,634\pm0,014$ & $16,862\pm0,018$ 
  \end{tabular}
  \label{tau:pendent_amplituds}
\end{Figura} 
\begin{Figura}
  \centering
  \includegraphics[width=1\linewidth]{../../graphs/practica_Ia/plots/reg_ampli.png}
  \captionof{figure}{\footnotesize{Regressió lineal del logaritme de les amplituds respecte a la posició. Coeficients de regressió obtinguts: $r^2$ = 0.999 tant per a la barra gran com per a la barar petita.}}
  \label{fig:reg_lin_amplituds}
\end{Figura}
Tenint en compte el comportament del desfasament i de les amplituds podem dir que les ones tèrmiques no són simètriques espaialment, decauen i pateixen un desfasament amb la distància.\\\\
Podem comprovar aquets resultats calculant el valor de $m$ a partir de l'Equació \ref{trobar_m}. Aplicant logaritme podem calcular $m$ per a qualsevol parell de posicions i qualsevol parell d'amplituds al mateix instant de temps. Escollim com a amplituds de referència les amplituds màximes trobades anteriorment. Calculem $m$ per a cada parell de màxims i fem la mitjana de tots per a cada barra. Aquest pas el repetim per a cada parell de posicions possible (veure Annex). Fent la mitjana de cada valor trobat acabem obtenint
\begin{equation*}
  \bar{m}_{\text{gran}} = (88,30\pm0,14)\cdot 10^{-3} \, \text{cm}^{-1}
\end{equation*}
per a la barra gran i
\begin{equation*}
  \bar{m}_{\text{petita}} = (86,343\pm0,037)\cdot 10^{-3} \, \text{cm}^{-1}
\end{equation*} 
per a la barra petita.\\\\
Comparant amb els pendent obtinguts a la Taula \ref{tau:pendent_amplituds} observem com, dins de la incertesa, els valors coincideixen.\\\\
Calculem ara el periode de les oscil·lacions. Per a fer-ho mesurem la distància entre els màxims trobats anteriorment, tot per a cada posició i cada barra. Trobem que, com calia esperar, aquests periodes d'oscil·lació són molt similars entre ells (veure Annex). Per tant, assumint que aquest periode és només un, donem el seu valor a partir de la mitjana de les dades trobades:
\begin{equation*}
  \bar{\tau} = (654,048\pm0,001) \, \text{s}
\end{equation*}
Sabent aquesta dada podem calcular el temps d'escalfament o de refredament mig senzillament dividint el periode a la meitat:
\begin{equation*}
  t_{\text{esc/ref}} = (325,024\pm0.001) \, \text{s}
\end{equation*}
Fent servir els màxims per a cada $x_i$ d'una mateixa barra podem calcular el desfasament entre dos punts. Tenint en compte l'Equació \eqref{increment_desfasament} podem calcular el valor de $h$ per a cada parell $x_i$. Fent la mitjana d'aquests valors (veure Annex) i apliacant-ho a cada barra trobem
\begin{equation*}
  \bar{h}_{\text{gran}} = (0,092\pm0,012) \, \text{cm}^{-1}
\end{equation*}
per a la barra gran i
\begin{equation*}
  \bar{h}_{\text{petita}} = (0,088\pm0,006) \, \text{cm}^{-1}
\end{equation*}
per a la barra petita.\\\\
Sabent tant els valors de $h$ i de $m$ per a cada barra podem calcular el valor de $\lambda$ segons l'Equació \ref{valor_lambda}. Fem servir el valor de la conductivitat tèrmica de l'alumini trobat a la literatura a la secció anterior i com a radis fem servir
\begin{equation*}
  r_{\text{gran}} = (2,6\pm0.1)\, \text{cm}
\end{equation*}
\begin{equation*}
  r_{\text{petita}} = (1,5\pm0.1)\, \text{cm}
\end{equation*}
per a la barra gran i la barra petita, respectivament. Aleshores, calculem els valors de $\lambda$ i obtenim
\begin{equation*}
  \lambda_{\text{gran}} = (-0.21\pm0.72)\cdot 10^{-2} \, \text{W}/\text{Kcm}^{2}
\end{equation*}
\begin{equation*}
  \lambda_{\text{petita}} = (-0.07\pm0.19)\cdot 10^{-2} \, \text{W}/\text{K}\text{cm}^{2}
\end{equation*}
Observem com el valor del coeficient de convecció pren un valor amb molta incertesa. També es pot notar com aquest coeficient creix amb l'àrea de la barra. Tot i això, no podem dir que els resultats coincideixen amb la realitat. Comparant amb resultats calculats amb les dades de la secció anterior aquests coeficients poden entrar dins del esperat, però només per a causa de tenir un error molt alt. Observem aquests coeficients a la Taula \ref{lambdas_alum_est}.
\begin{Figura}
  \centering
  \captionof{table}{\footnotesize{Coeficients de convecció per a una barra d'alumini gran i petita (\(r_g=2.6\,\text{cm}\) i \(r_p =1.5\,\text{cm}\)) calculats amb les dades de la secció de l'estat estacionari.}}
  \begin{tabular}{c|c}
    Barra & $\lambda$ $(\cdot 10^{-2})$ [W/cmK]\\ \hline \hline
    Gran & $0,291\pm0,042$\\
    Petita & $0,168\pm0,026$ 
  \end{tabular}
\end{Figura}
\section{Conclusions}
Al llarg de la pràctica s'ha posat de manifest com les diferents equacions de transport de calor defineixen el comportament del nostre sistema. Tenint en compte aquestes aportacions de cada equació, hem obtingut dues solucions per a diferents sistemes termodinàmics.\\\\
Per al primer sistema trobem una solució estacionària. Amb les mesures experimentals hem observat com el comportament de la temperatura a cada barra és coherent amb la solució del nostre model. Aquestes dades ens han permès trobar relacions entre els coeficients de conducció tèrmica dels materials amb que treballem. Comparant amb dades de la literatura, els nostres resultats són coherents per a la barra de ferro i per a la barra d'alumini. Tot i això, per a la barra de llautó els resultats difereixen, exposant que possiblement el model falla en algún punt, i aquest error pot provenir de diverses fonts. Dos possibles fonts (a part de l'error experimental) són considerar el coeficient de convecció constant per als tres metalls, una altra, considerar que el coeficient de conducció tèrmica no depèn de la temperatura. Els nostres resultats experimentals també han permès observar la poca aportació de potència que prové de la radiació.\\\\
Per al segón sistema trobem una solució permanent. Les característiques d'aquest sistema fan que a cada barra es generi una ona tèrmica que es propaga longitudinalment. Hem observat que la solució del model és coherent amb les nostres dades experimentals, podent així calcular diferents paràmetres que defineixen la ona tèrmica, com el període de les oscil·lacions o el temps de refredament o escalfament. Hem observat que, encara que el comportament sigui oscil·latori, encara se segueix complint el decaiment exponencial de la temperatura envers la distància, de forma anàloga a com ho feia el sistema anterior. No només això, s'ha comprovat experimentalment també que existeix un desfasament de la ona que depèn de la distància. Aquest desfasament és coherent amb la solució del model.
\section{Annex}
\subsection{Càlcul d'incerteses}
Per a calcular la incertesa dels valors mesurats directament al laboratori s'ha fet servir la desviació estàndar de la mostra obtinguda de cada valor mesurat. Juntament amb la incertesa instrumental, calculem la incertesa de cada mesura com
    \begin{equation*}
        u(x_i) = \sqrt{ s_{\text{est}}^2(x_i) + s_{\text{ins}}^2}
    \end{equation*}
    on $s_{\text{est}}$ és la desviació estàndar, $s_{\text{ins}}$ és la incertesa instrumental i $x_i$ és el valor mesurat.\\\\
    Per a calcular la incertesa de magnituds derivades de les variables mesurades, s'ha fet servir la propagació d'incerteses, calculada com
    \begin{equation*}
        u^2(y) = \sum_{i = 1}^{n} \left( \frac{\partial y}{\partial{x_i}} u(x_i)\right)^2
    \end{equation*}
    on $y$ és una funció tal que $y = y(x_1,...,x_n)$.\\\\
    Per a calcular els errors de les regressions s'ha fet servir la funció \texttt{curve\_fit} de la llibreria \texttt{scipy.optimize} de Python. Per a calcular els coeficients de correlació $r^2$ s'ha aplicat el mètode de mínims quadrats també amb Python.
    \subsection{Dades experimentals}
    Les següents taules contenen diversos valors mesurats i calculats al llarg de la pràctica:
    \begin{Figura}
      \centering
      \captionof{table}{\footnotesize{Temperatures mesurades per al sistema en l'estat estacionari per a l'alumini.}}
      \begin{tabular}{c|c}
        $x$ [cm] & $T_{\text{Al}}$ ($\pm$ 0,1) [$^\circ$C]\\\hline \hline
        110 & 29,3 \\
        100 & 30,0 \\
        90 &  30,3 \\
        80 &  31,6\\
        70 &  33,7\\
        60 &  40,6\\
        50 &  46,0\\
        40 &  53,1\\
        30 &  63,7\\
        20 &  77,6\\
        10 &  95,5\\
        0 &  119,0\\
      \end{tabular}
    \end{Figura}
    \begin{Figura}
      \centering
      \captionof{table}{\footnotesize{Temperatures mesurades per al sistema en l'estat estacionari per al llautó.}}
      \begin{tabular}{c|c}
        $x$ [cm] & $T_{\text{Cu-Zn}}$ ($\pm$ 0,1) [$^\circ$C]\\\hline \hline
        110 & 26,6\\
        100 & 27,0\\
        90 &  27,6\\
        80 &  28,7\\
        70 &  30,5\\
        60 &  32,5\\
        50 &  35,8\\
        40 &  39,9\\
        30 &  44,8\\
        20 &  55,2\\
        10 &  72,8\\
        0 &  108,9\\
      \end{tabular}
    \end{Figura}
    \begin{Figura}
      \centering
      \captionof{table}{\footnotesize{Temperatures mesurades per al sistema en l'estat estacionari per al ferro.}}
      \begin{tabular}{c|c}
        $x$ [cm] & $T_{\text{Fe}}$ ($\pm$ 0,1) [$^\circ$C]\\\hline \hline
        110 &  24,9\\
        100 &  25,3\\
        90 &  25,3\\
        80 &  25,8\\
        70 &  26,2\\
        60 &  26,7\\
        50 &  27,9\\
        40 &  30,2\\
        30 &  34,7\\
        20 &  44,3\\
        10 &  65,0\\
        0 &  120,6\\
      \end{tabular}
    \end{Figura}
    \begin{Figura}
      \centering
      \captionof{table}{\footnotesize{Mesura dels diàmetres de cada barra.}}
      \begin{tabular}{c|ccc}
        Materials & \multicolumn{3}{c}{$\varnothing$($\pm$0,001) [cm]}\\ \hline \hline
        Al & 3,035 & 3,065 & 3,045\\
        Cu-Zn & 3,02 & 3,025 & 3,065\\
        Fe & 3,145 & 3,285 & 3,01  
      \end{tabular}
    \end{Figura}
    \begin{Figura}
      \centering
      \captionof{table}{\footnotesize{Mesures de la temperatura ambient del laboratori.}}
      \begin{tabular}{c}
          $T_{\text{amb}}$ ($\pm$0,1) [$^\circ$C] \\ \hline\hline
          24,7\\
          24,9\\
          24,6 
      \end{tabular}
      \label{T_ambient}
    \end{Figura}
    \begin{Figura}
      \centering
      \captionof{table}{\footnotesize{Periodes d'oscil·lació calculats per a cada posició i cada barra.}}
      \begin{tabular}{c|c}
        $x_i$ [cm] & $\tau$ ($\pm$0,001) [s]\\ \hline\hline
        \multicolumn{2}{c}{Barra gran}\\\hline
        0 & 654,281 \\
        10 & 655,686 \\
        15 & 657,090 \\\hline
        \multicolumn{2}{c}{Barra petita}\\\hline
        0 & 654,280 \\
        10 & 652,879 \\
        20 & 650,069
      \end{tabular}
    \end{Figura}
    \begin{Figura}
      \centering
      \captionof{table}{\footnotesize{Valors de $m$ per a diferents increments $\Delta x_i$ per a la barra gran}}
      \begin{tabular}{c|c}
        $\Delta x_i$ [cm]  & $m_{\text{gran}}$ ($\cdot 10^{-3}$) [cm$^{-1}$]\\ \hline\hline
        10 - 0 & $88,899\pm0,040$ \\
        15 - 10 & $88,90\pm0,14$ \\
        15 - 0 & $87,100\pm0,017$
      \end{tabular}
    \end{Figura}
    \begin{Figura}
      \centering
      \captionof{table}{\footnotesize{Valors de $m$ per a diferents increments $\Delta x_i$ per a la barra petita}}
      \begin{tabular}{c|c}
        $\Delta x_i$ [cm]  & $m_{\text{petita}}$ ($\cdot 10^{-3}$) [cm$^{-1}$]\\ \hline\hline
        10 - 0 & $86,591\pm0,037$ \\
        20 - 10 & $86,096\pm0,037$ \\
        20 - 0 & $86,3431\pm0,0093$
      \end{tabular}
    \end{Figura}
    \begin{Figura}
      \centering
      \captionof{table}{\footnotesize{Valors de $h$ per a diferents increments $\Delta x_i$ per a la barra gran}}
      \begin{tabular}{c|c}
        $\Delta x_i$ [cm]  & $h_{\text{gran}}$ [cm$^{-1}$]\\ \hline\hline
        10 - 0 & $0,109\pm0,007$ \\
        15 - 10 & $0,072\pm0,013$ \\
        15 - 0 & $0,096\pm0,004$
      \end{tabular}
    \end{Figura}
    \begin{Figura}
      \centering
      \captionof{table}{\footnotesize{Valors de $h$ per a diferents increments $\Delta x_i$ per a la barra petita}}
      \begin{tabular}{c|c}
        $\Delta x_i$ [cm]  & $h_{\text{gran}}$ [cm$^{-1}$]\\ \hline\hline
        10 - 0 & $0,082\pm0,005$ \\
        20 - 10 & $0,095\pm0,006$ \\
        20 - 0 & $0,088\pm0,003$
      \end{tabular}
    \end{Figura}
\end{multicols}
\bibliographystyle{plain} 
\bibliography{referencias_Ia.bib}
\end{document}
